% Abhakseite „Das kann ich“ – Zone nur im Gesamt (\mitzone)
%% TODO Ich-kann-Sätze: Platzhalter wie in den Aufgabendateien
\begin{abhakseite}
\ifdefined\mitzone
\abhakgruppe{Kennst du schon}
\abhak{1}{Skalarprodukt in Koordinaten berechnen, Beträge bilden}
\abhak{2}{Normalenvektoren aus Koordinatengleichungen ablesen und aus Spannvektoren bestimmen}
\abhak{3}{Richtungsvektoren aus Geraden und Kanten bilden}
\abhak{4}{Kosinus und Tangens am Taschenrechner auswerten und die Umkehrfunktion anwenden (Grad-Modus)}
\abhak{5}{Neben-, Komplement- und Innenwinkel unterscheiden und ergänzen}
\abhak{6}{Kreisumfang und Anteile, Prozentrechnung}
\abhak{7}{Neben-, Komplement- und Innenwinkel unterscheiden und ergänzen – Fehler finden}
\abhak{8}{Neben-, Komplement- und Innenwinkel unterscheiden und ergänzen – selbst rechnen}
\fi
\abhakgruppe{Einheit 1 · Das Skalarprodukt als Objekt}
\abhak{9}{Das Skalarprodukt als Objekt}
\abhak{10}{Das Skalarprodukt als Objekt – Fehler finden, Begründen}
\abhakgruppe{Einheit 2 · Winkel zwischen Vektoren, Kanten und Geraden}
\abhak{11}{Der Formelwinkel oder sein Nachbar?}
\abhak{12}{Winkel zwischen Vektoren, Kanten und Geraden}
\abhak{13}{Winkel zwischen Vektoren, Kanten und Geraden – Fehler finden, Begründen, Anwendung}
\abhakgruppe{Einheit 3 · Neigungswinkel von Ebenen}
\abhak{14}{Neigungswinkel von Ebenen}
\abhak{15}{Neigungswinkel von Ebenen – Fehler finden, Begründen, Anwendung}
\abhakgruppe{Einheit 4 · Winkel von Geraden gegen Ebenen und Bogenmaße}
\abhak{16}{Winkel gegen Ebenen und Bogenmaße}
\abhak{17}{Winkel gegen Ebenen und Bogenmaße}
\abhak{18}{Winkel gegen Ebenen und Bogenmaße – Fehler finden, Begründen, Anwendung}
\end{abhakseite}
